
% language=us runpath=texruns:manuals/ontarget

\startcomponent ontarget-inclusion

\environment ontarget-style

\startchapter[title={Including \PDF\ files}]

Rendering glyphs in a \PDF\ happens based on information in the page stream. In
such a stream we find font triggers that use identifiers like \type {F1} and
afterwards placement operators inject shapes by referring to an index in the
font, say hexadecimal \type {0003}, and that index can refer to any shape.
The identifier is resolved via the \type {Font} entry in the \type
{Resources} dictionary of
the page:

\starttyping
5 0 obj
<<
    /Contents 3 0 R
    /Resources << /Font << /F1 1 0 R >> /ProcSet 2 0 R >>
    ...
    /Type /Page
>>
endobj
\stoptyping

Following the \type {F1} reference we end up at:

\starttyping
1 0 obj
<<
    /BaseFont /TNTUFI+LMRoman10-Regular
    /DescendantFonts [ 11 0 R ]
    /Encoding /Identity-H
    /Subtype /Type0
    /ToUnicode 14 0 R
    /Type /Font
>>
endobj
\stoptyping

and when we then descend into the first entry in the \typ {DescendantFonts} array we come
to:

\starttyping
11 0 obj
<<
    /BaseFont /TNTUFI+LMRoman10-Regular
    /CIDSystemInfo << /Ordering (Identity) /Registry (Adobe) /Supplement 0 >>
    /FontDescriptor 9 0 R
    /LMTXRegistry 8 0 R
    /Subtype /CIDFontType0
    /Type /Font
    /W 10 0 R
>>
endobj
\stoptyping

The \type {W} entry value, rather short relative to the other keys, refers to an array object
that holds the widths of the glyphs so that the viewer knows how much to advance;
the \typ {FontDescriptor} points to the shapes; and \typ {LMTXRegistry} will be
discussed later. These examples come from typesetting:

\starttyping
\starttext
    \startTEXpage
        test
    \stopTEXpage
\stoptext
\stoptyping

With \LUATEX\ we see this in the page stream:

\starttyping
BT
/F1 11.955168 Tf
1 0 0 1 0 0.11949 Tm [<0069003200620069>]TJ
ET
\stoptyping

and also this in the mapping from index to \UNICODE, which is object
14, defined by the \typ
{ToUnicode} value shown above (I added the characters as comments):

\starttyping
3 beginbfchar
<0032> <0065> % e
<0062> <0073> % s
<0069> <0074> % t
endbfchar
\stoptyping

When we use \LUAMETATEX\ instead we get:

\starttyping
BT
/F1 10 Tf
1.195517 0 0 1.195517 0 0.065717 Tm [<0001000200030001>] TJ
ET
\stoptyping

and:

\starttyping
3 beginbfchar
<0001> <0074> % t
<0002> <0065> % e
<0003> <0073> % s
endbfchar
\stoptyping

Thus, where \LUATEX\ uses the original index in the font (not to be confused with
the character's \UNICODE\ value, if it has one at all), in \LUAMETATEX, or more accurately with the
\CONTEXT\ backend, we start at one and number upwards. This gives smaller files.

The only way to find out what an index is actually referring to is to
consult the abovementioned \type {ToUnicode} vector in the font resource because
there we map from index to \UNICODE. That information is used when you search in
a \PDF\ file or cut|-|and|-|paste from it. Because glyphs can be unrelated to \UNICODE,
and because multiple glyphs can share the same \UNICODE\ slot, the index is what
makes a glyph unique.

When a (page from) a \PDF\ file is included in a document \LUATEX\ will copy the
relevant objects to the main file. Of course the page itself is copied (with the
page stream making up the content). Copying is also driven by the \type
{Resources} key in the page dictionary. In addition to the \type {Font}
list we've seen above, there can also be an \type {XObject} array and its entries need to be copied
as well. This copying is recursive because the resources themselves can point to
objects with resources. It is quite normal in a \PDF\ file to share resources.
The rendered glyphs, for instance, come from fonts that contain the shape
definitions and basically these are the same, independent of scaling.

The index can be the original glyph index in the font but, because we subset, it
can also be a different one, depending on what gets included. This is
illustrated in the example above. By default the \LUATEX\ engine just copies and
doesn't worry about what gets copied. However, in \MKIV\ we can load some code
that plugs into the \LUATEX\ backend and is thereby capable of merging fonts from
the included \PDF\ (page) with one used in the document. This process is driven
by setting the \type {compact} key in \typ {\externalfigure} to the value \type
{yes}. Here we assume that the references to glyphs in the page stream are the
original (or equivalent) indices but this is not guaranteed to be true. We can check a
little by comparing the \UNICODE\ mapping as well as doing a visual check
afterwards, but neither are robust. It still works ok as long we use exactly
the same font; essentially, we check the name and when it matches we force the glyph
into the current file and use the font reference of main document for the
embedded reference instead.

In \LUAMETATEX\ we do it differently and there are reasons for this. First of
all, we have different numbering in the main file and inserted file, so we cannot
use the indices directly. In addition, we have to also take care of \TYPETHREE\
fonts that refer to fonts that we merge (we use these fonts in, for instance, math
delimiters). Finally we cannot simply look at the name because we can have a
variable font instance that has different axis properties. So we have to be more
clever: we need to parse the content streams of the page, XObjects and charprocs
to find out what glyphs are referenced and replace indices when applicable. In
addition we consult some extra information that is included when \CONTEXT\ did
typeset the (to be embedded) file. That information contains a stream index to
original index mapping, and also has a variable font recipe if needed.
There is also some additional information so that we at least check if we have the same
font.

\starttabulate[|l|l|l|c|c|c|]
\FL
\NC             \NC       \NC               \BC native            \BC \type{compact=no} \BC \type{compact=yes} \NC \NR
\ML
\BC \LUATEX     \BC \MKIV \BC compressed    \NC  839 KB / .20 sec \NC 731 KB / 0.20 sec \NC 231 KB / 0.22 sec  \NC \NR
\NC             \NC       \BC decompressed  \NC 1784 KB / .20 sec \NC 943 KB / 0.22 sec \NC 426 KB / 0.23 sec  \NC \NR
\ML
\BC \LUAMETATEX \BC \MKXL \BC compressed    \NC                   \NC 544 KB / 0.18 sec \NC 147 KB / 0.24 sec  \NC \NR
\NC             \NC       \BC decompressed  \NC                   \NC 783 KB / 0.16 sec \NC 552 KB / 0.19 sec  \NC \NR
\LL
\stoptabulate

we compare three alternatives. The native inclusion in \LUATEX\
leaves the work to the engine. When the plugin is loaded, we will use the
\LUA|-|based inclusion code which is a bit more clever in sharing objects. In the
\LUAMETATEX\ variant we also copy objects into the main file but there we have no
plugin and sharing happens anyway. Here with \typ {compact=yes} we also merge
fonts but this time based on parsing the streams. This parsing is more demanding
and bumps runtime but is also more rewarding in terms of file size. For
completeness we show the results with and without \PDF\ object stream (zip)
compression. The need to decompress and compress also has some
impact on performance.

In case the slightly slower inclusion in \LUAMETATEX\ bothers you, there might be
some comfort in knowing that the 20 files accumulate to 564 \KB\ and a fresh run
takes 49 seconds. When we use \LUATEX\ the file size total bumps to 748 \KB\ and
the initial runtime goes up to 94 seconds. So in the end the \LUAMETATEX|-|based
variant is more efficient.

So, in \MKIV\ there are two reasons for having the plugin. The first is that by
sharing common objects we can, for instance, include many pages from the same
document with little overhead. For this, compact doesn't need to be active.
However when for instance we include many documents we can see that merging fonts
does pay off handsomely. In \MKXL\ we already have the first benefit (sharing) by
default and here we can also do better by merging fonts. Because that merging is
more aggressive you see better numbers in the table for \MKXL.

A practical usage scenario is making manuals where we process
examples (using \CONTEXT\
buffers) in independent runs so that they are independent from the main document.
Of course there is only a gain if these examples share fonts with the main document
or with each other. Here is the test case:

\starttyping
\startbuffer[common]
    \usebodyfont  [dejavu]
    \usebodyfont  [lucida]
    \usebodyfont  [bonum]
    \setupbodyfont[modern]
    \setupalign[tolerant,stretch]
\stopbuffer
\stoptyping

We load four different font sets in a common buffer but also use them in the
main document:

\starttyping
\getbuffer[common]
\stoptyping

We define two additional buffers that each create a document with two pages. We
use different languages because otherwise there is little to merge, as
the sample texts use the regular Latin script:

\starttyping
\startbuffer[demo-1]
    \start
        \switchtobodyfont[dejavu]\samplefile{ward}
    \stop \blank
    \start
        \switchtobodyfont[lucida]\samplefile{davis}
    \stop \page
    \start
        \switchtobodyfont[bonum] \samplefile{knuth}
    \stop \blank
    \start
        \switchtobodyfont[modern]\samplefile{tufte}
    \stop \blank
\stopbuffer

\startbuffer[demo-2]
    \start
        \switchtobodyfont[dejavu]\mainlanguage[es]
        \samplefile{cervantes-es.tex}
    \stop \blank
    \start
        \switchtobodyfont[lucida]\mainlanguage[sv]
        \samplefile{alfredsson-sv.tex}
    \stop \page
    \start
        \switchtobodyfont[bonum] \mainlanguage[de]
        \samplefile{aesop-de.tex}
    \stop \blank
    \start
        \switchtobodyfont[modern]\mainlanguage[cz]
        \samplefile{komensky-cz}
    \stop \blank
\stopbuffer
\stoptyping

Optionally we enable compact inclusion:

\starttyping
% \setupexternalfigures[compact=yes]
\stoptyping

Here is the main document:

\starttyping
\starttext
    \dorecurse{10}{
        \startTEXpage[offset=1ex]
            \start \switchtobodyfont[dejavu]\samplefile{ward}  \stop \blank
            \start \switchtobodyfont[lucida]\samplefile{davis} \stop \blank
            \start \switchtobodyfont[bonum] \samplefile{knuth} \stop \blank
            \start \switchtobodyfont[modern]\samplefile{tufte} \stop \blank
            \setbuffer[#1]#1\endbuffer % #1 is the iterator
            \hbox\bgroup
                \typesetbuffer[common,demo-1,#1][width=10cm,page=1]
                \typesetbuffer[common,demo-1,#1][width=10cm,page=2]
            \egroup
            \blank
            \hbox\bgroup % these are processed in sperate runs
                \typesetbuffer[common,demo-2,#1][width=10cm,page=1]
                \typesetbuffer[common,demo-2,#1][width=10cm,page=2]
            \egroup
        \stopTEXpage
    }
\stoptext
\stoptyping

In order to get ten times two unique documents to be included, we smuggle an
extra buffer into the subsidiary runs. We need to do this because otherwise the hashes
of the content of these sub|-|documents are the same and we'd end up with only two
documents and successive inclusions would share these. Of course the first run
with fresh buffers will take more runtime because the sub|-|documents need to be
processed (twice in order to get multi|-|pass activities resolved).

There are a few pitfalls. First of all we have to make sure that we only merge
references to the same font. This is often no problem as long as we don't update
fonts with different shapes in the same slots but we can assume that the version
number is different then. For the application we have in mind, buffered
sub|-|runs or inclusion in related documents that get processed in a short time span, we
are probably ok. We can make the check more tolerant or more clever, and
might do that in the future. On the average the inclusion is already rather
efficient when a few pages from a few documents are used.

A second pitfall is
that when we improve the \CONTEXT\ (font) backend we can have better
shapes or more precise metrics but because metrics are unlikely to change much in
the glyph programs we're probably okay. Even mixing the more efficient so|-|called
compact font mode with normal font mode (not to be confused with compact
inclusion) should work out well enough. In case of doubt: trust your eyes or just
regenerate the documents involved in the inclusion.

Finally it is worth mentioning
that there is a noticeable overhead but if becomes necessary I can optimize the
handling of the stream a bit by replacing the more general stream parser by a
dedicated one for this purpose.

Let's stress one thing again. Because shared font usage will never be (guaranteed)
watertight you do need to check visually. A bad merge will immediately show up by
the included image rendering with garbled text. There are some extra safeguards
in the \MKXL\ approach that are absent in the \MKIV\ solution which is why the
latter is considered an experiment and not loaded by default. I could spend time on
it but as we moved on to \LMTX\ (the \MKXL\ \LUAMETATEX\ combination) it makes
little sense.

Does the story end here? Not entirely. Occasional validation
requirements have a side effect that
some users have to fix old \PDF\ files to suit demands. Let's mention a few issues
users run into:

\startitemize[packed]
    \startitem
        Embedded so|-|called \TYPEZERO\ fonts can lack a \type {/CIDSet} and|/|or
        \type {CIDToGIDMap} entry that needs to be added.
    \stopitem
    \startitem
        A document can refer to external files that are not embedded. Normally these
        are in \type {WinAnsi} encoding and page stream indices match encoding indices
        so we can smuggle these files into the document.
    \stopitem
    \startitem
        Font resources can be embedded multiple times with different subsets, likely
        per page. Different names are used, which complicates matters, but it makes
        sense to try to merge them.
    \stopitem
    \startitem
        Fonts with the same name but in \TYPEONE\ as well as \TRUETYPE\ format, both
        using the original indices, occur in the same document so they can be merged.
    \stopitem
\stopitemize

We can deal with this quite well if we have the original fonts available.
Failures to do this well immediately show up so we can again trust our eyes. In an
automatic large scale fix operation we can built in some safeguards.

Then, as we're fixing included pages anyway, we may as well try to conform to
standard even better, for instance:

\startitemize[packed]
    \startitem
        Instead of gray scales \CMYK\ and \RGB\ colors are used, or one of these
        color spaces is not handled right and we need to remap colors. It can be
        a side effect of lazy programming in producer software.
    \stopitem
    \startitem
        Extended graphic states are used, for instance for transparencies while
        there is no transparency actually being used. These can be side effects of
        producers just emitting as much as they can.
    \stopitem
    \startitem
        Irrelevant grouping can be applied to pages and \type {/XObjects}. In
        fact, when a validator complains we can just as well get rid of them.
    \stopitem
\stopitemize

For such things we need to fix the \type{/Resources} as well as the page content streams
so it adds a little more overhead but when we just convert documents it is a
one|-|time effort so a few more milliseconds won't be a big burden.

But discussing these details doesn't really fit in this font discussion so we
end by mentioning that we have a framework in place for plugging in fixers of
any kind. The approach is to configure what standard to use and what fixes to
apply. Only time will tell if this is sufficient. Most users probably never have
to worry about it anyway.

{\em Many thanks to Karl Berry for copy|-|editing this text!}

\stopchapter

\stopcomponent

%
% % engine=luatex
%
% % \nopdfcompression
%
% % \ifcase\contextlmtxmode
% %     \registerctxluafile{grph-bmp}{}
% %     \registerctxluafile{grph-chk}{}
% % \fi
%
% \startbuffer[common]
%     \usebodyfont  [dejavu]
%     \usebodyfont  [lucida]
%     \usebodyfont  [bonum]
%     \setupbodyfont[modern]
%     \setupalign[tolerant,stretch]
% \stopbuffer
%
% % \enabletrackers[graphics.fonts]
% \setupexternalfigures[compact=yes]
%
%
% \getbuffer[common]
%
% \startbuffer[demo-1]
%     \start \switchtobodyfont[dejavu]\samplefile{ward}  \stop \blank
%     \start \switchtobodyfont[lucida]\samplefile{davis} \stop \blank
%     \page
%     \start \switchtobodyfont[bonum] \samplefile{knuth} \stop \blank
%     \start \switchtobodyfont[modern]\samplefile{tufte} \stop \blank
% \stopbuffer
%
% \startbuffer[demo-2]
%     \start \switchtobodyfont[dejavu]\mainlanguage[es]\samplefile{cervantes-es.tex}  \stop \blank
%     \start \switchtobodyfont[lucida]\mainlanguage[sv]\samplefile{alfredsson-sv.tex} \stop \blank
%     \page
%     \start \switchtobodyfont[bonum] \mainlanguage[de]\samplefile{aesop-de.tex} \stop \blank
%     \start \switchtobodyfont[modern]\mainlanguage[cz]\samplefile{komensky-cz}  \stop \blank
% \stopbuffer
%
% \starttext
%
% \testfeatureonce{1} {
%     \dorecurse{10}{
%         \startTEXpage[offset=1ex]
%             \start \switchtobodyfont[dejavu]\samplefile{ward}  \stop \blank
%             \start \switchtobodyfont[lucida]\samplefile{davis} \stop \blank
%             \start \switchtobodyfont[bonum] \samplefile{knuth} \stop \blank
%             \start \switchtobodyfont[modern]\samplefile{tufte} \stop \blank
%             \normalexpanded{\setbuffer[\recurselevel]\recurselevel\endbuffer}
%             \hbox\bgroup
%                 \typesetbuffer[common,demo-1,\recurselevel][width=10cm,page=1]
%                 \typesetbuffer[common,demo-1,\recurselevel][width=10cm,page=2]
%             \egroup
%             \blank
%             \hbox\bgroup
%                 \typesetbuffer[common,demo-2,\recurselevel][width=10cm,page=1]
%                 \typesetbuffer[common,demo-2,\recurselevel][width=10cm,page=2]
%             \egroup
%         \stopTEXpage
%     }
% }
%
% \startTEXpage
%     \strut \elapsedtime
% \stopTEXpage
%
% \stoptext

